\section{数据结构}
\subsection{单调队列}
$O(n)$ 滑动窗口最值问题
\begin{lstlisting}[caption={}]
//求最大构建递减的单调队列 
for(int i = 0;i < n;i++){
	while(!Q.empty() && a[i] >= a[Q.back()]) Q.pop_back();
	Q.push_back(i);
    if(i < k - 1) continue;
	while(!Q.empty() && Q.front() <= i - k) Q.pop_front();
}
\end{lstlisting}


\subsection{单调栈}
$O(n)$ 对于每个位置找到其左边第一个比它大的位置
\begin{lstlisting}[caption={}]
std::vector<int> stk;
for(int i = 1;i <= n;i++){
	while(!stk.empty() && a[stk.back()][0] <= a[i][0]) stk.pop_back();
	if(!stk.empty()) ans[stk.back()] += a[i][1];
	stk.push_back(i);
}
\end{lstlisting}

\subsection{并查集}
\begin{lstlisting}[caption={}]
struct DSU{
    vector<ll> p;
    DSU ();
    DSU(int n){init(n);};
    void init(int n){
        p.resize(n + 2);
        for(int i = 1;i <= n;i++)
            p[i] = i;
    }
    ll find(ll x){
        if(p[x] != x) p[x] = find(p[x]);
        return p[x];
    }
    void merge(ll x,ll y){
        x = find(x);
        y = find(y);
        p[y] = x;
    }
};
\end{lstlisting}

\subsection{ST表}
ST表是用于解决可重复贡献问题的数据结构,如$max(x,y),gcd(a,b)$等。
\begin{lstlisting}[caption={}]
#include<bits/stdc++.h>
#define all(x) (x).begin(),(x).end()
using namespace std;
typedef long long ll;
const int N = 1e5 + 10;

//P3865 【模板】ST 表
//静态区间最值
struct ST{
	vector<vector<ll>> st;
	ll N,op;
	ST(vector<ll>& s,int op){
		N = s.size();
		st.resize(N + 1,vector<ll>(20,-1e18));
		for(int i = 1;i <= N;i++){
			st[i][0] = s[i];
		}
		this->op = op;
		build_st(op);
	}

	ll opt1(const ll a,const ll b){return max(a,b);}
	ll opt2(const ll a,const ll b){return min(a,b);}

	ll query(ll l,ll r){
		ll k = log2(r - l + 1); 
		if(op == 1) return opt1(st[l][k],st[r - (1 << k) + 1][k]); 
		return opt2(st[l][k],st[r - (1 << k) + 1][k]); 
	}
	void build_st(int op){
		for(int j = 1;j <= 20;j++){
			for(int i = 1;i + (1 << j) - 1 <= N;i++){
				if(op == 1) st[i][j] = opt1(st[i][j - 1],st[i + (1 << (j - 1))][j - 1]);
				else st[i][j] = opt2(st[i][j - 1],st[i + (1 << (j - 1))][j - 1]);
			}
		}
	}
};
\end{lstlisting}

\subsection{线段树}
\subsubsection{区间求和}
线段树可以在$O(\log N)$的时间复杂度内实现单点修改、区间修改、区间查询（区间求和，求区间最大值，求区间最小值）等操作。
\begin{lstlisting}[caption={}]
struct SegmentTree{
    ll n;
    struct node{
        ll date,tag;
    };
    vector<node> tr;
    SegmentTree(vector<ll>& nums){
        n = nums.size() - 1;
        tr.resize(n * 4 + 10);
        build(1,n,1,nums);
    }
    void push_up(ll p){
        tr[p].date = tr[p << 1].date + tr[(p << 1) | 1].date;
    }
    void f(ll p,ll l,ll r,ll x){
        tr[p].tag += x;
        tr[p].date += x * (r - l + 1);
    }
    void push_down(ll p,ll l,ll r){
        ll mid = (l + r) >> 1LL;
        f(p << 1,l,mid,tr[p].tag);
        f(p << 1 | 1,mid + 1,r,tr[p].tag);
        tr[p].tag = 0;
    }
    void build(ll l,ll r,ll p,vector<ll>& nums){
        if(l == r){
            tr[p].date = nums[l];
            return;
        }
        ll m = l + ((r - l) >> 1);
        build(l,m,p << 1,nums),build(m + 1,r,(p << 1) | 1,nums);
        push_up(p);
    }
    void update(ll l,ll r,ll s,ll t,ll p,ll c){
        if(l <= s && t <= r){
            tr[p].date += (t - s + 1) * c;
            tr[p].tag += c;
            return;
        }
        push_down(p,s,t);
        ll m = s + ((t - s) >> 1);
        if(l <= m) update(l,r,s,m,p << 1,c);
        if(r > m) update(l,r,m + 1,t,(p << 1) | 1,c);
        push_up(p);
    }
    ll query(ll l,ll r,ll s,ll t,ll p){
        if(l <= s && t <= r) return tr[p].date;
        push_down(p,s,t);
        ll m = s + ((t - s) >> 1);
        ll sum = 0;
        if(l <= m) sum = query(l,r,s,m,p << 1);
        if(r > m) sum += query(l,r,m + 1,t,(p << 1) | 1);
        return sum;
    }
    void update(ll l,ll r,ll x){
        update(l,r,1,n,1,x);
    }
    ll query(ll l,ll r){
        return query(l,r,1,n,1);
    }
};
\end{lstlisting}

\subsubsection{区间加乘}
\begin{lstlisting}[caption={}]
struct SegmentTree{
    ll n,mod = 998244353;
    struct node{
        ll date = 0,add = 0,mu = 1;
    };
    vector<node> tr;
    SegmentTree(vector<ll>& nums){
        n = nums.size() - 1;
        tr.resize(n * 4 + 10);
        build(1,n,1,nums);
    }
    void push_up(ll p){
        tr[p].date = (tr[p << 1].date + tr[(p << 1) | 1].date) % mod;
    }
    void f(ll p,ll l,ll r,ll x,ll k){
        tr[p].date = (tr[p].date * k + x * (r - l + 1)) % mod;
        tr[p].mu = (tr[p].mu * k) % mod;
        tr[p].add = (tr[p].add * k % mod + x) % mod;
    }
    void push_down(ll p,ll l,ll r){
        ll mid = (l + r) >> 1LL;
        f(p << 1,l,mid,tr[p].add,tr[p].mu);
        f(p << 1 | 1,mid + 1,r,tr[p].add,tr[p].mu);
        tr[p].add = 0,tr[p].mu = 1;
    }
    void build(ll l,ll r,ll p,vector<ll>& nums){
        if(l == r){
            tr[p].date = nums[l];
            return;
        }
        ll m = l + ((r - l) >> 1);
        build(l,m,p << 1,nums),build(m + 1,r,(p << 1) | 1,nums);
        push_up(p);
    }
    //x is add,k is mul
    void update(ll l,ll r,ll s,ll t,ll p,ll x,ll k){
        if(l <= s && t <= r){
            tr[p].date = (tr[p].date * k % mod + (t - s + 1) * x) % mod;
            tr[p].mu = (tr[p].mu * k) % mod;
            tr[p].add = (tr[p].add * k % mod + x) % mod;
            return;
        }
        push_down(p,s,t);
        ll m = s + ((t - s) >> 1);
        if(l <= m) update(l,r,s,m,p << 1,x,k);
        if(r > m) update(l,r,m + 1,t,(p << 1) | 1,x,k);
        push_up(p);
    }
    ll query(ll l,ll r,ll s,ll t,ll p){
        if(l <= s && t <= r) return tr[p].date;
        push_down(p,s,t);
        ll m = s + ((t - s) >> 1);
        ll sum = 0;
        if(l <= m) sum = query(l,r,s,m,p << 1);
        if(r > m) sum = (sum + query(l,r,m + 1,t,(p << 1) | 1)) % mod;
        return sum;
    }
    void RangeMul(ll l,ll r,ll x){
        update(l,r,1,n,1,0,x);
    }
    void RangeAdd(ll l,ll r,ll x){
        update(l,r,1,n,1,x,1);
    }
    ll query(ll l,ll r){
        return query(l,r,1,n,1);
    }
};
\end{lstlisting}

\subsubsection{Info + Tag}
\begin{lstlisting}[caption={}]
template<class Info,class Tag>
struct SegmentTree{
    i64 N;
    std::vector<Info> info;
    std::vector<Tag> tag;
    SegmentTree():N(0){};
    SegmentTree(int _n,Info _v = Info()){init(_n,_v);}
    template<class T>
    SegmentTree(std::vector<T> _init){init(_init);}
    void init(int _n,Info _v = Info()){init(std::vector<Info>(_n,_v));}

    template<class T>
    void init(std::vector<T> _init){
        N = _init.size();
        info.assign(4 << std::__lg(N),Info());
        tag.assign(4 << std::__lg(N),Tag());
        auto build = [&](auto&& self,int p,int l,int r) -> void{
            if(l == r){
                info[p] = _init[l];
                return;
            }
            int m = (l + r) >> 1;
            self(self,p << 1,l,m);
            self(self,(p << 1) | 1,m + 1,r);
            push_up(p);
        };
        build(build,1,1,N);
    }
    void push_up(int p){
        info[p] = info[p << 1] + info[(p << 1) | 1];
    }
    void apply(int p,const Tag &v){
        info[p].apply(v);
        tag[p].apply(v);
    }
    void push_down(int p){
        apply(p << 1,tag[p]);
        apply((p << 1) | 1,tag[p]);
        tag[p] = Tag();
    }
    void modify(int p,int l,int r,int x,int y,const Tag &v){
        if(l > y || r < x) return;
        if(x <= l && r <= y){
            apply(p,v);
            return;
        }
        push_down(p);
        i64 m = (l + r) >> 1;
        modify(p << 1,l,m,x,y,v);
        modify((p << 1) | 1,m + 1,r,x,y,v);
        push_up(p);
    }
    void modify(int l, int r, const Tag &v) {
        return modify(1,1,N,l,r,v);
    }
    Info query(int p,int l,int r,int x,int y){
        if(l > y || r < x) return Info();
        if(x <= l && r <= y) return info[p];
        push_down(p);
        i64 m = (l + r) >> 1;
        return query(p << 1,l,m,x,y) + query((p << 1) | 1,m + 1,r,x,y);
    }
    Info query(int l,int r){
        return query(1,1,N,l,r);
    }
};
struct Tag{
    i64 x = 0;
    void apply(const Tag &t) & {
        x += t.x;
    }
};
struct Info{
    i64 x = 0,len = 1;
    void apply(const Tag &t) & {
        x += t.x * len;
    }
    
    Info operator+(const Info &other) const{
        return Info{x + other.x,len + other.len};
    }
};
\end{lstlisting}

\subsection{树状数组}
树状数组是一种支持 单点修改 和 区间查询 的，代码量小的数据结构。普通树状数组维护的信息及运算要满足 结合律 且 可差分，如加法（和）、乘法（积）、异或等。
\begin{lstlisting}[caption={}]
struct BIT{
    vector<i64> tr;
    BIT():tr(N){};
    void update(i64 p,i64 x){
        for(int i = p;i <= 1e5;i += i & -i)
            tr[i] += x;
    }
    i64 query(i64 p){
        i64 res = 0;
        for(int i = p;i >= 1;i -= i & -i)
            res += tr[i];
        return res;
    }
};
\end{lstlisting}
前驱后继扩展（常规+第 k 小值查询+元素排名查询+元素前驱后继查询）

注意，被查询的值都应该小于等于$N$，否则会越界；如果离散化不可使用，则需要使用平衡树替代。
\begin{lstlisting}[caption={}]
struct BIT {
    int n;
    vector<int> w;
    BIT(int n) : n(n), w(n + 1) {}
    void add(int x, int v) {
        for (; x <= n; x += x & -x)
            w[x] += v;
    }
    int kth(int x) { // 查找第 k 小的值
        int ans = 0;
        for (int i = __lg(n); i >= 0; i--) {
            int val = ans + (1 << i);
            if (val < n && w[val] < x) {
                x -= w[val];
                ans = val;
            }
        }
        return ans + 1;
    }
    int get(int x) { // 查找 x 的排名
        int ans = 1;
        for(x--; x; x -= x & -x)
            ans += w[x];
        return ans;
    }
    int pre(int x) { return kth(get(x) - 1); } // 求小于 x 的最大值（前驱）
    int suf(int x) { return kth(get(x + 1)); } // 求大于 x 的最小值（后继）
};
\end{lstlisting}

\subsection{李超线段树}
加入一个一次函数，定义域为$[l,r]$,给定$k$,求定义域包含$k$的所有一次函数中,
在$x=k$处取值最大的那个,如果有多个函数取值相同,选编号最小的。
\begin{lstlisting}[caption={}]
struct Lichao{
    #define ls (p << 1)
    #define rs ((p << 1) | 1)
    static constexpr double eps = 1e-8;
    struct line{
        double k,b;
        int id;      // 部分题需要记录线段编号
        line(double k = 0,double b = 0,int id = 0):k(k),b(b),id(id){};
        double calc(int x){ return k * x + b;} // 计算线段在pos位置的纵坐标
        double cross(const line& T){ return (T.b - b) / (k - T.k);} // 求两条线段的交点横坐标
    };
    struct node{
        int l,r;   // 区间左右端点
        line L;    // 最优线段
        bool flag; // 记录区间内是否有最优线段
        node(){};
        node(int l,int r,line L,bool flag):l(l),r(r),L(L),flag(flag){};
    };
    std::vector<node> tr;
    Lichao(){}
    Lichao(int N):tr(N << 2){
        build(1,1,N);
    };
    void build(int p,int l,int r){
        tr[p] = {l,r,line(),false};
        if(l == r) return;
        int mid = (l + r) >> 1;
        build(ls,l,mid),build(rs,mid + 1,r);
    }
    //O(log^2 n)
    void update(int p,int L,int R,line ln){
        int l = tr[p].l,r = tr[p].r;
        if(L <= l && r <= R){ //如果插入线段的定义域覆盖整个区间
            double lp = tr[p].L.calc(l),rp = tr[p].L.calc(r);
            double lq = ln.calc(l),rq = ln.calc(r);

            if(!tr[p].flag) tr[p].L = ln,tr[p].flag = true;
            else if(lq - lp > eps && rq - rp > eps) tr[p].L = ln;
            else if(lq - lp > eps || rq - rp > eps){
                int mid = (l + r) >> 1;
                if(ln.calc(mid) - tr[p].L.calc(mid) > eps) std::swap(tr[p].L,ln);
                if(ln.calc(l) - tr[p].L.calc(l) > eps) update(ls,L,R,ln);
                else update(rs,L,R,ln);
            }
        }
        else{
            int mid = (l + r) >> 1;
            if(L <= mid) update(ls,L,R,ln);
            if(R > mid) update(rs,L,R,ln);
        }
    }
    //O(log n)
    std::pair<double,int> query(int p,int pos){ //返回最大值和线段编号
        int l = tr[p].l,r = tr[p].r;
        double ans = tr[p].L.calc(pos);
        int id = tr[p].L.id;
        if(l == r) return {ans,id};
        int mid = (l + r) >> 1;
        std::pair<double,int> tq;
        if(pos <= mid) tq = query(ls,pos);
        else tq = query(rs,pos);

        if(tq.first > ans || (fabs(tq.first - ans) < eps && tq.second < id))
            ans = tq.first,id = tq.second;
        return {ans,id};
    }
    std::pair<double,int> query(int x){
        return query(1,x);
    }
    void update(i64 x0,i64 y0,i64 x1,i64 y1,i64 id){
        line ln;
        if(x0 == x1) ln = {0,std::max(y1,y0),id};
        else{
            double tk = (double)(y1 - y0) / (x1 - x0);
            double tb = (double)(y0 - tk * x0);
            ln = {tk,tb,id};
        }
        update(1,x0,x1,ln);
    }
};
\end{lstlisting}

\subsection{动态开点李超线段树}
\begin{lstlisting}[caption={}]
struct Lichao{
    static constexpr double eps = 1e-8;
    struct Line{
        double k,b;
        int id;
        Line(double k = 0,double b = 0,int id = 0):k(k),b(b),id(id){};
        double calc(int x){return k * x + b;}
        double cross(const Line& T){return (T.b - b) / (k - T.k);}
    };
    struct Node{
        int l,r,ls,rs;
        Line L;
        bool flag;
        Node(){};
        Node(int l,int r):l(l),r(r),ls(0),rs(0),flag(false){};
    };
    std::vector<Node> tr;
    int cnt;
    int L_range,R_range;
    Lichao(int L, int R):L_range(L),R_range(R),cnt(1){
        tr.push_back(Node());
        tr.push_back(Node(L,R));
    }
    void new_node(int& p,int l,int r){
        p = ++cnt;
        tr.push_back(Node(l,r));
    }
    void update(int p,int L,int R,Line ln){
        int l = tr[p].l, r = tr[p].r;
        int mid = (l + r) >> 1;
        if(L <= l && r <= R){
            if(!tr[p].flag){
                tr[p].L = ln;
                tr[p].flag = true;
                return;
            }
            double lp = tr[p].L.calc(l),rp = tr[p].L.calc(r);
            double lq = ln.calc(l),rq = ln.calc(r);
            if(!tr[p].flag) tr[p].L = ln,tr[p].flag = true;
            else if(lq - lp > eps && rq - rp > eps) tr[p].L = ln;
            else if(lq - lp > eps || rq - rp > eps){
                int mid = (l + r) >> 1;
                if(ln.calc(mid) - tr[p].L.calc(mid) > eps) std::swap(tr[p].L,ln);
                if(ln.calc(l) - tr[p].L.calc(l) > eps){
                    if(!tr[p].ls) new_node(tr[p].ls,l,mid);
                    update(tr[p].ls,L,R,ln);
                }
                else{
                    if(!tr[p].rs) new_node(tr[p].rs,mid + 1,r);
                    update(tr[p].rs,L,R,ln);
                }
            }
        }
        else{
            if(L <= mid){
                if(!tr[p].ls) new_node(tr[p].ls,l,mid);
                update(tr[p].ls,L,R,ln);
            }
            if(R > mid){
                if(!tr[p].rs) new_node(tr[p].rs,mid + 1,r);
                update(tr[p].rs,L,R,ln);
            }
        }
    }
    std::pair<double,int> query(int p,int pos){
        int l = tr[p].l,r = tr[p].r;
        double ans = tr[p].L.calc(pos);
        int id = tr[p].L.id;
        if(l == r) return {ans,id};
        int mid = (l + r) >> 1;

        std::pair<double,int> tq;
        if(pos <= mid && tr[p].ls) tq = query(tr[p].ls, pos);
        else if(pos > mid && tr[p].rs) tq = query(tr[p].rs, pos);
        else tq = {-1e18,0};
        
        if(tq.first > ans || (fabs(tq.first - ans) < eps && tq.second < id))
            ans = tq.first,id = tq.second;
        return {ans,id};
    }
    std::pair<double,int> query(int x){
        return query(1,x);
    }
    void update(i64 x0,i64 y0,i64 x1,i64 y1,i64 id){
        Line ln;
        if(x0 == x1) ln = {0,std::max(y0,y1),id};
        else{
            double tk = (double)(y1 - y0) / (x1 - x0);
            double tb = y0 - tk * x0;
            ln = {tk, tb,id};
        }
        update(1,x0,x1,ln);
    }
};
\end{lstlisting}

\subsection{珂朵莉树}
珂朵莉树的适用范围是有区间赋值操作且数据随机的题目。在随机数据下，珂朵莉树可以达到$O(n\log \log n)$的复杂度。
\begin{lstlisting}[caption={}]
struct node{
	ll l, r;
	mutable ll v;
	node(ll l, ll r, ll v) : l(l), r(r), v(v) {}
	bool operator<(const node &o) const { return l < o.l; }
};
struct ODT{
    set<node> tr;
    ODT(){};
    auto split(ll p){
        auto t = tr.lower_bound({p,0,0});
        if(t != tr.end() && t->l == p) return t;
        t--;
        auto [l,r,v] = *t;
        tr.erase(t);
        tr.insert({l,p - 1,v});
        return tr.insert({p,r,v}).first;
    }
    void assign(ll l,ll r,ll v){
        auto end = split(r + 1), begin = split(l);
		tr.erase(begin, end);
		tr.insert(node(l,r,v));
    }
	void add(ll l,ll r,ll v){
		auto end = split(r + 1);
		for (auto it = split(l); it != end; it++)
			it->v += v;
	}
	ll kth(ll l,ll r,ll k){
		auto end = split(r + 1);
		vector<PII> v;
		for (auto it = split(l); it != end; it++)
			v.push_back({it->v, it->r - it->l + 1});
		sort(v.begin(), v.end());
		for (int i = 0; i < v.size(); i++){
			k -= v[i].second;
			if (k <= 0)
				return v[i].first;
		}
	}
	ll sum_of_pow(ll l, ll r, ll x, ll y){
		ll tot = 0;
		auto end = split(r + 1);
		for (auto it = split(l); it != end; it++)
			tot = (tot + qpow(it->v, x, y) * (it->r - it->l + 1)) % y;
		return tot;
	}
};
\end{lstlisting}

\subsection{主席树(可持久化线段树)}
以 $\mathcal O(N\log N)$ 的时间复杂度建树、查询、修改。
\begin{lstlisting}[caption={}]
struct PresidentTree {
    static constexpr int N = 2e5 + 10;
    int cntNodes, root[N];
    struct node {
        int l, r;
        int cnt;
    } tr[4 * N + 17 * N];
    void modify(int &u, int v, int l, int r, int x) {
        u = ++cntNodes;
        tr[u] = tr[v];
        tr[u].cnt++;
        if (l == r) return;
        int mid = (l + r) / 2;
        if (x <= mid)
            modify(tr[u].l, tr[v].l, l, mid, x);
        else
            modify(tr[u].r, tr[v].r, mid + 1, r, x);
    }
    int kth(int u, int v, int l, int r, int k) {
        if (l == r) return l;
        int res = tr[tr[v].l].cnt - tr[tr[u].l].cnt;
        int mid = (l + r) / 2;
        if (k <= res)
            return kth(tr[u].l, tr[v].l, l, mid, k);
        else
            return kth(tr[u].r, tr[v].r, mid + 1, r, k - res);
    }
};
\end{lstlisting}



\subsection{笛卡尔树}
\begin{lstlisting}[caption={}]
template<typename T>
struct CarTree{
	struct node{
		int l,r;
		node():l(0),r(0){}
	};
	vector<node> tr;
	CarTree(){};
	CarTree(vector<T>& a){
		stack<T> stk;
		tr.resize(a.size());
		for(int i = 1;i < a.size();i++){
			int last = 0;
			while(!stk.empty() && a[stk.top()] > a[i]){
				last = stk.top();
				stk.pop();
			}
			if(last != 0) tr[i].l = last;
			if(!stk.empty()) tr[stk.top()].r = i;
			stk.push(i);
		}
	}
};
\end{lstlisting}

\subsection{pbds 扩展库实现平衡二叉树}
记得加上相应的头文件，同时需要注意定义时的参数，一般只需要修改第三个参数：即定义的是大根堆还是小根堆。
\begin{lstlisting}[caption={}]
empty() / size()
insert(x) // 插入元素x
erase(x) // 删除元素/迭代器x
order_of_key(x) // 返回元素x的排名
find_by_order(x) // 返回排名为x的元素迭代器
lower_bound(x) / upper_bound(x) // 返回迭代器
join(Tree) // 将Tree树的全部元素并入当前的树
split(x, Tree) // 将大于x的元素放入Tree树
\end{lstlisting}
\begin{lstlisting}[caption={}]
#include <ext/pb_ds/assoc_container.hpp>
using namespace __gnu_pbds;
using V = pair<int, int>;
tree<V, null_type, less<V>, rb_tree_tag, tree_order_statistics_node_update> ver;
map<int, int> dic;

int n; cin >> n;
for (int i = 1, op, x; i <= n; i++) {
    cin >> op >> x;
    if (op == 1) { // 插入一个元素x，允许重复
        ver.insert({x, ++dic[x]});
    } else if (op == 2) { // 删除元素x，若有重复，则任意删除一个
        ver.erase({x, dic[x]--});
    } else if (op == 3) { // 查询元素x的排名（排名定义为比当前数小的数的个数+1）
        cout << ver.order_of_key({x, 1}) + 1 << endl;
    } else if (op == 4) { // 查询排名为x的元素
        cout << ver.find_by_order(--x)->first << endl;
    } else if (op == 5) { // 查询元素x的前驱
        int idx = ver.order_of_key({x, 1}) - 1; // 无论x存不存在，idx都代表x的位置，需要-1
        cout << ver.find_by_order(idx)->first << endl;
    } else if (op == 6) { // 查询元素x的后继
        int idx = ver.order_of_key( {x, dic[x]}); // 如果x不存在，那么idx就是x的后继
        if (ver.find({x, 1}) != ver.end()) idx++; // 如果x存在，那么idx是x的位置，需要+1
        cout << ver.find_by_order(idx)->first << endl;
    }
}
\end{lstlisting}

\subsection{vector模拟平衡树}
\begin{lstlisting}[caption={}]
int main(){
    std::ios::sync_with_stdio(0),std::cin.tie(0),std::cout.tie(0);
    i64 n;
    std::cin >> n;
    std::vector<i64> a;
    while(n--){
        i64 opt,x;
        std::cin >> opt >> x;
        //插入O(log n)
        if(opt == 1)      a.insert(lower_bound(a.begin(),a.end(),x),x);
        //删除O(log n)
        else if(opt == 2) a.erase(lower_bound(a.begin(),a.end(),x));
        //查询排名O(log n)
        else if(opt == 3) std::cout << lower_bound(a.begin(),a.end(),x) - a.begin() + 1 << "\n";
        //查询排名对应数字O(1)
        else if(opt == 4) std::cout << a[x - 1] << "\n";
        //前驱O(log n) 
        else if(opt == 5) std::cout << *(lower_bound(a.begin(),a.end(),x) - 1) << "\n";
        //后继O(log n)
        else              std::cout << *upper_bound(a.begin(),a.end(),x) << "\n";
    }
    return 0;
}
\end{lstlisting}